% !Mode:: "TeX:UTF-8"
% 开源 starter 模板，保留天津大学本科毕业论文常用排版设置。

%%%%%%%%%% Package %%%%%%%%%%%%
% 使用 ctexbook + LuaLaTeX/XeLaTeX 时不应再加载 CJK 宏包（会干扰字体设置）
\def\atemp{xelatex}\def\btemp{lualatex}
\ifx\atemp\usewhat\else\ifx\btemp\usewhat\else
  \usepackage{CJK}
\fi\fi
\usepackage{lmodern}
\usepackage[T1]{fontenc}
\usepackage{graphicx}                       % 支持插图处理
\usepackage[a4paper,top= 27.5true mm,bottom=25.4 true mm,left=35.7 true mm,right=27.7 true mm, headheight=15true mm,footskip=8true mm]{geometry}
                                            % 支持版面尺寸设置
% \usepackage[squaren]{SIunits}             % 旧宏包，Overleaf 上易与现有命令冲突
\usepackage{titlesec}                       % 控制标题的宏包
\usepackage{titletoc}                       % 控制目录的宏包
\usepackage{fancyhdr}                       % fancyhdr宏包 支持页眉和页脚的相关定义
\def\atemp{xelatex}\def\btemp{lualatex}
\ifx\atemp\usewhat\else\ifx\btemp\usewhat\else
  \def\atemp{xelatex}\def\btemp{lualatex}
\ifx\atemp\usewhat\else\ifx\btemp\usewhat\else
  \usepackage{CJKpunct}                       % 精细调整中文的标点符号
\fi\fi
\fi\fi
\usepackage{color}                          % 支持彩色
\usepackage{amsmath}                        % AMSLaTeX宏包 用来排出更加漂亮的公式
\usepackage{amssymb}                        % 数学符号生成命令
\usepackage[below]{placeins}    %允许上一个section的浮动图形出现在下一个section的开始部分,还提供\FloatBarrier命令,使所有未处理的浮动图形立即被处理
\usepackage{multirow}                       % 使用Multirow宏包，使得表格可以合并多个row格
\usepackage{booktabs}                       % 表格，横的粗线；\specialrule{1pt}{0pt}{0pt}
\usepackage{makecell}                       % 支持单元格合并
\usepackage{longtable}                      % 支持跨页的表格。
\usepackage{tabularx}                       % 自动设置表格的列宽
% \usepackage{ccaption}                       % 支持子图的中文标题
% \usepackage{subfig}                         
\usepackage{booktabs}                       % 支持行内多表格
\usepackage{subcaption}                     % 支持子表
\captionsetup[subtable]{labelformat=empty}  % 子表不标序号
\usepackage[numbers]{natbib}  % 不压缩区间，避免出现 [1-6] 这类过宽泛引用
\usepackage{float}                         % 支持 [H] 固定浮动体位置
\usepackage{enumitem}                       % 使用enumitem宏包,改变列表项的格式
\usepackage{calc}                           % 长度可以用+ - * / 进行计算
\usepackage{txfonts}                        % 字体宏包
\usepackage{bm}                             % 处理数学公式中的黑斜体的宏包
\usepackage[amsmath,thmmarks,hyperref]{ntheorem}  % 定理类环境宏包，其中 amsmath 选项用来兼容 AMS LaTeX 的宏包
\def\atemp{xelatex}\def\btemp{lualatex}
\ifx\atemp\usewhat\else\ifx\btemp\usewhat\else
  \def\atemp{xelatex}\def\btemp{lualatex}
\ifx\atemp\usewhat\else\ifx\btemp\usewhat\else
  \usepackage{CJKnumb}                        % 提供将阿拉伯数字转换成中文数字的命令
\fi\fi
\fi\fi
\usepackage{indentfirst}                    % 首行缩进宏包
\def\atemp{xelatex}\def\btemp{lualatex}
\ifx\atemp\usewhat\else\ifx\btemp\usewhat\else
  \def\atemp{xelatex}\def\btemp{lualatex}
\ifx\atemp\usewhat\else\ifx\btemp\usewhat\else
  \usepackage{CJKutf8}                        % 用在UTF8编码环境下，它可以自动调用CJK，同时针对UTF8编码作了设置
\fi\fi
\fi\fi
\newcommand{\tabincell}[2]{\begin{tabular}{@{}#1@{}}#2\end{tabular}}
\usepackage{xcolor}
\usepackage{tikz}
% 支持表格中叉与对勾
\usepackage{pifont}
\DeclareRobustCommand{\cmark}{\textcolor{green!60!black}{\ding{51}}} % ✓
\DeclareRobustCommand{\xmark}{\textcolor{red}{\ding{55}}} % ✗
\usepackage[table]{xcolor} % 提供 \columncolor
\usepackage{threeparttable}
\usepackage{adjustbox}
\usepackage{xstring}
\usepackage{pgf}          % \pgfmath... 计算
\definecolor{NatureBlue}{RGB}{34,94,168}
\definecolor{BTUp}{RGB}{200,60,60}     % 增加：红
\definecolor{BTDown}{RGB}{60,110,190}  % 减少：蓝
\makeatletter
\def\bt@dmax{1} % 默认值，防止未设置时报错
\newcommand{\btsetdmax}[1]{\def\bt@dmax{#1}}
\newcommand{\btcell}[2]{%
  \begingroup
  \def\bt@val{#1}%
  \def\bt@delta{#2}%
  \expandafter\btcell@parse\bt@delta\relax
  \endgroup
}
\def\btcell@parse(#1#2)\relax{%
  \pgfmathsetmacro{\bt@abs}{abs(#2)}%
  \pgfmathsetmacro{\bt@ratio}{min(1, \bt@abs/(\bt@dmax))}%
  \pgfmathtruncatemacro{\bt@int}{round(6 + 26*\bt@ratio)}%
  \def\bt@base{NatureBlue}%
  \ifx#1+%
    \def\bt@base{BTUp}%
  \else\ifx#1-%
    \def\bt@base{BTDown}%
  \fi\fi
  \edef\bt@colorspec{\bt@base!\bt@int}%
  \expandafter\cellcolor\expandafter{\bt@colorspec}%
  \bt@val{\scriptsize\textcolor{gray}{\,\bt@delta}}%
}
\makeatother
%支持代码环境
\usepackage{listings}
\lstset{numbers=left,
language=[ANSI]{C},
numberstyle=\tiny,
extendedchars=false,
showstringspaces=false,
breakatwhitespace=false,
breaklines=true,
captionpos=b,
keywordstyle=\color{blue!70},
commentstyle=\color{red!50!green!50!blue!50},
frame=shadowbox,
rulesepcolor=\color{red!20!green!20!blue!20}
}
%支持算法环境
\usepackage[boxed,ruled,lined,linesnumbered]{algorithm2e}
\renewcommand{\algorithmcfname}{算法}
\SetAlFnt{\it}
% 算法按章编号，与图、表分章编号习惯一致
\makeatletter
\@addtoreset{algocf}{chapter}
\makeatother
\renewcommand{\thealgocf}{\thechapter-\arabic{algocf}}
\SetKwInput{KwIn}{输入}
\SetKwInput{KwOut}{输出}
\SetKw{KwTo}{至}
\SetKw{KwRet}{返回}
\SetKw{KwCont}{continue}

\usepackage{array}
\newcommand{\PreserveBackslash}[1]{\let\temp=\\#1\let\\=\temp}
\newcolumntype{C}[1]{>{\PreserveBackslash\centering}p{#1}}
\newcolumntype{R}[1]{>{\PreserveBackslash\raggedleft}p{#1}}
\newcolumntype{L}[1]{>{\PreserveBackslash\raggedright}p{#1}}

% XeLaTeX / LuaLaTeX 下这两段配置一致，直接统一加载即可
\usepackage{pdfpages}
\usepackage[unicode,
            pdfstartview=FitH,
            bookmarksnumbered=true,
            bookmarksopen=true,
            plainpages=false,
            pdfpagelabels=true,
            colorlinks=false,
            pdfborder={0 0 0},
            citecolor=blue,
            linkcolor=red,
            anchorcolor=green,
            urlcolor=blue,
            breaklinks=true
            ]{hyperref}
\pdfstringdefDisableCommands{%
  \def\\{ }%
}
